\lysh{34186}{4}

\begin{alist}
\item Η εξίσωση $x^2 - 2x + \lambda(2 - \lambda) = 0$ έχει ρίζες $x_1$ και $x_2$ με άθροισμα $S = x_1 + x_2 = -\dfrac{-2}{1} = 2$
και γινόμενο $P = x_1 x_2 = \dfrac{\lambda(2 - \lambda)}{1} = \lambda(2 - \lambda)$.
\begin{rlist}
\item Η περίμετρος $\Pi$ του ορθογωνίου είναι: 
\[\Pi = 2x_1 + 2x_2 = 2(x_1 + x_2) = 2 \cdot 2 = 4.\]

\item Το εμβαδόν $E$ του ορθογωνίου ως συνάρτηση του $\lambda$ είναι:
\[E = x_1 x_2 = \lambda(2 - \lambda), \ \ \lambda \in (0, 2).\]
\end{rlist}

\item Έχουμε
\begin{align*}
E \le 1 \ &\Leftrightarrow \\
\lambda(2 - \lambda) \le 1 \ &\Leftrightarrow \\
\lambda^2 - 2\lambda + 1 \ge 0 \ &\Leftrightarrow
\end{align*}
$(\lambda - 1)^2 \ge 0$, που ισχύει για κάθε $\lambda \in (0, 2)$.

\item Το εμβαδόν $E$ του ορθογωνίου γίνεται μέγιστο, δηλαδή ίσο με $1$, αν και μόνο αν
\[E = 1 \stackrel{(\beta)}{\Leftrightarrow} (\lambda - 1)^2 = 0 \Leftrightarrow \lambda = 1.\]

Για $\lambda = 1$, η περίμετρος του ορθογωνίου είναι $4$ και το εμβαδόν του $1$, οπότε το ορθογώνιο
γίνεται τετράγωνο πλευράς $x_1 = x_2 = 1$.
\end{alist}

